% abkuerzungen.tex
% Abkuerzungsverzeichnis der Bachelorarbeit.
% Wird in main.tex mit % abkuerzungen.tex
% Abkuerzungsverzeichnis der Bachelorarbeit.
% Wird in main.tex mit % abkuerzungen.tex
% Abkuerzungsverzeichnis der Bachelorarbeit.
% Wird in main.tex mit % abkuerzungen.tex
% Abkuerzungsverzeichnis der Bachelorarbeit.
% Wird in main.tex mit \input{abkuerzungen} eingebunden.
%
% Verwendung im Text:
%   \ac{API}   -- beim ersten Vorkommen: "Application Programming Interface (API)"
%                 danach automatisch nur noch: "API"
%   \acl{API}  -- immer die Langform: "Application Programming Interface"
%   \acs{API}  -- immer die Kurzform: "API"

\begin{acronym}[WYSIWYG]  % Laengste Abkuerzung fuer Einrueckung angeben
  \acro{API}[API]{Application Programming Interface}
  \acro{BA}[BA]{Bachelorarbeit}
  \acro{CI}[CI]{Continuous Integration}
  \acro{CD}[CD]{Continuous Deployment}
  \acro{GUI}[GUI]{Graphical User Interface}
  \acro{HTTP}[HTTP]{Hypertext Transfer Protocol}
  \acro{REST}[REST]{Representational State Transfer}
  \acro{SPA}[SPA]{Single-Page-Applikation}
  \acro{UI}[UI]{User Interface}
  \acro{URL}[URL]{Uniform Resource Locator}
  \acro{UX}[UX]{User Experience}
  \acro{WYSIWYG}[WYSIWYG]{What You See Is What You Get}
\end{acronym}
 eingebunden.
%
% Verwendung im Text:
%   \ac{API}   -- beim ersten Vorkommen: "Application Programming Interface (API)"
%                 danach automatisch nur noch: "API"
%   \acl{API}  -- immer die Langform: "Application Programming Interface"
%   \acs{API}  -- immer die Kurzform: "API"

\begin{acronym}[WYSIWYG]  % Laengste Abkuerzung fuer Einrueckung angeben
  \acro{API}[API]{Application Programming Interface}
  \acro{BA}[BA]{Bachelorarbeit}
  \acro{CI}[CI]{Continuous Integration}
  \acro{CD}[CD]{Continuous Deployment}
  \acro{GUI}[GUI]{Graphical User Interface}
  \acro{HTTP}[HTTP]{Hypertext Transfer Protocol}
  \acro{REST}[REST]{Representational State Transfer}
  \acro{SPA}[SPA]{Single-Page-Applikation}
  \acro{UI}[UI]{User Interface}
  \acro{URL}[URL]{Uniform Resource Locator}
  \acro{UX}[UX]{User Experience}
  \acro{WYSIWYG}[WYSIWYG]{What You See Is What You Get}
\end{acronym}
 eingebunden.
%
% Verwendung im Text:
%   \ac{API}   -- beim ersten Vorkommen: "Application Programming Interface (API)"
%                 danach automatisch nur noch: "API"
%   \acl{API}  -- immer die Langform: "Application Programming Interface"
%   \acs{API}  -- immer die Kurzform: "API"

\begin{acronym}[WYSIWYG]  % Laengste Abkuerzung fuer Einrueckung angeben
  \acro{API}[API]{Application Programming Interface}
  \acro{BA}[BA]{Bachelorarbeit}
  \acro{CI}[CI]{Continuous Integration}
  \acro{CD}[CD]{Continuous Deployment}
  \acro{GUI}[GUI]{Graphical User Interface}
  \acro{HTTP}[HTTP]{Hypertext Transfer Protocol}
  \acro{REST}[REST]{Representational State Transfer}
  \acro{SPA}[SPA]{Single-Page-Applikation}
  \acro{UI}[UI]{User Interface}
  \acro{URL}[URL]{Uniform Resource Locator}
  \acro{UX}[UX]{User Experience}
  \acro{WYSIWYG}[WYSIWYG]{What You See Is What You Get}
\end{acronym}
 eingebunden.
%
% Verwendung im Text:
%   \ac{API}   -- beim ersten Vorkommen: "Application Programming Interface (API)"
%                 danach automatisch nur noch: "API"
%   \acl{API}  -- immer die Langform: "Application Programming Interface"
%   \acs{API}  -- immer die Kurzform: "API"

\begin{acronym}[WYSIWYG]  % Laengste Abkuerzung fuer Einrueckung angeben
  \acro{API}[API]{Application Programming Interface}
  \acro{BA}[BA]{Bachelorarbeit}
  \acro{CI}[CI]{Continuous Integration}
  \acro{CD}[CD]{Continuous Deployment}
  \acro{GUI}[GUI]{Graphical User Interface}
  \acro{HTTP}[HTTP]{Hypertext Transfer Protocol}
  \acro{REST}[REST]{Representational State Transfer}
  \acro{SPA}[SPA]{Single-Page-Applikation}
  \acro{UI}[UI]{User Interface}
  \acro{URL}[URL]{Uniform Resource Locator}
  \acro{UX}[UX]{User Experience}
  \acro{WYSIWYG}[WYSIWYG]{What You See Is What You Get}
\end{acronym}
